% Response to Reviewer 2
% Same text as rebuttal/reviewer_2.md (1491 characters as posted).
\section*{Response to Reviewer 2}

Thank you for a thoughtful review and for valuing our training-free interface and honest reporting.

\textbf{W1:} VID and K400 are the standard benchmarks, but their short clips prevent a long-video test, as we will state. Over videos of a few hundred frames, one early memory kept the full saving, a fixed share of the host's budget per frame. With rebuilds every $N$ frames, a long video is a chain of such segments, so we expect results to hold; dynamic vs.\ static scenes remain open.

\textbf{W2:} Our main claim: PSM layers onto Eventful without retraining ($\gamma = 0$ is the host), cutting 60.7 to 46.6 GFLOPs/frame. We agree the key test is the host at equal cost: a budget-matched Eventful recovers $\sim$1.5 mAP (p.~5), so we will narrow our claims; combining memory and change signals is next.

\textbf{W3:} Sorry for the inconsistency: no threshold is used; $\tau$ (p.~3) is a leftover from an earlier version. Algorithm 1's fraction $\gamma$ skips the $\gamma k$ proposed tokens most similar to a memory rebuilt every $N$ frames (unused in our runs); skipped tokens keep the host's K/V (l.~14). We will rewrite the method text.

\textbf{W4:} We report latency to show it stays near the host and 14\% below dense, not to claim a speed-up; even the host turns 65\% fewer FLOPs into only 16\% lower latency. We acknowledge the larger K400 losses at higher budgets and will narrow those claims.

\textbf{W5:} Both tasks follow prior work and test two readouts, not one; the readout explanation stays a hypothesis, as we will say.

Thank you again!
